% Consolidated 15-transport from 350eef03dcd4ca0e99eacbc61d97caa887d8c65a; original TeX fragments preserved.

% === BEGIN 15-transport-220-229.tex blob 159513f0c7ed913d408c1f0c023350880352784d ===
% Irodov 用户PDF页97--98，原书印刷页88--89；第2.7节220--229题。
% 使用硬球碰撞平均自由程 lambda=(sqrt2 pi d^2 n)^-1。
\thermotransportproblem{220}{88}
求气体分子的自由程长度：
(a) 超过平均自由程 $\lambda$ 的分数；
(b) 位于 $\lambda$ 与 $2\lambda$ 之间的分数。
\exsolution
硬球模型下分子从最近一次碰撞后，
在路程 $s$ 内未碰撞的存活概率为
$P(L>s)=e^{-s/\lambda}$。
(a) 直接代入 $s=\lambda$：
\[
\boxed{P(L>\lambda)=e^{-1}\simeq0.368}.
\]
(b) 区间概率等于两次存活概率之差：
\[
\boxed{P(\lambda\le L\le2\lambda)
=e^{-1}-e^{-2}\simeq0.233}.
\]
\excheck 自由程的平均值是 $\lambda$，
但其分布不是所有分子都飞过恰好 $\lambda$，
而是具有指数长尾。

\thermotransportproblem{221}{88}
一束狭窄分子束射入充满低压气体的容器。
沿传播方向经过距离 $\Delta\ell$ 后，
分子束强度降低为原来的 $1/\eta$。
求束流分子的平均自由程 $\lambda$。
\exsolution
未发生碰撞的分子数按
$I(\ell)=I_0e^{-\ell/\lambda}$ 衰减。
由
$I(\Delta\ell)/I_0=1/\eta$：
\[
e^{-\Delta\ell/\lambda}=\eta^{-1}.
\]
两边取自然对数：
$-\Delta\ell/\lambda=-\ln\eta$。
因此
\[
\boxed{\lambda=\frac{\Delta\ell}{\ln\eta}},
\quad\eta>1.
\]
\excheck 这里测到的是\emph{未碰撞的直束强度}，
要求散射分子不再被计入狭束探测器。

\thermotransportproblem{222}{88}
若气体分子在短时间
$\mathrm dt$ 内发生碰撞的概率为
$\alpha\,\mathrm dt$，其中碰撞率
$\alpha$ 为常数，求：
(a) 在时间 $t$ 内完全没有碰撞的概率；
(b) 两次相继碰撞间的平均等待时间。
\exsolution
设无碰撞存活概率为 $P_0(t)$。
在下一微时段，
$P_0(t+\mathrm dt)=P_0(t)(1-\alpha\mathrm dt)$，
故
$\mathrm dP_0/\mathrm dt=-\alpha P_0$。
用初态 $P_0(0)=1$，
得
\[
\boxed{P_0(t)=e^{-\alpha t}}.
\]
等待时间的概率密度是
$f(t)=-P_0'(t)=\alpha e^{-\alpha t}$，
平均值：
\[
\boxed{\bar\tau=\int_0^\infty
t\alpha e^{-\alpha t}\,\mathrm dt=\frac1\alpha}.
\]
\excheck $\alpha$ 的单位是 $\mathrm{s^{-1}}$，
而 $\bar\tau$ 的单位是秒。

\thermotransportproblem{223}{88}
求氮气分子的平均自由程 $\lambda$
及平均碰撞间隔 $\tau$：
(a) 标准状况；
(b) 温度 $0^\circ\mathrm C$，
压强 $p=1.0\,\mathrm{nPa}$ 的极高真空。
\exsolution
硬球模型中
\[
\lambda=\frac1{\sqrt2\,\pi d^2n}
=\frac{k_BT}{\sqrt2\,\pi d^2p},
\quad
\bar v=\sqrt{\frac{8RT}{\pi M}},
\quad \tau=\frac\lambda{\bar v}.
\]
取氮气分子有效直径
$d\simeq0.37\,\mathrm{nm}$、
$M\simeq0.028\,\mathrm{kg/mol}$、
$T\simeq273\,\mathrm K$。
(a) 标准压强
$p_0\simeq1.013\times10^5\,\mathrm{Pa}$，
得到数量级
\[
\boxed{\lambda_0\simeq0.06\,\mathrm{\mu m}},
\qquad
\boxed{\tau_0\simeq0.13\,\mathrm{ns}}.
\]
(b) 温度不变，自由程与压强成反比：
\[
\frac{\lambda_b}{\lambda_0}=\frac{p_0}{p_b}
\simeq1.0\times10^{14}.
\]
因此
\[
\boxed{\lambda_b\simeq6\times10^6\,\mathrm m
=6\,\mathrm{Mm}},\qquad
\boxed{\tau_b\simeq1.4\times10^4\,\mathrm s
\simeq3.8\,\mathrm h}.
\]
\excheck 本版答案经 PDF 文本提取后显示“6 mm”，
但与 $\tau=3.8\,\mathrm h$ 及
自由程反比压强的关系不相容；
应理解为 $6\,\mathrm{Mm}$（百万米）量级。
稀薄气体中该自由程远超通常容器尺寸，
真实分子大多先撞壁。

\thermotransportproblem{224}{88}
标准状况下，
氮气分子的平均自由程是
气体分子之间典型平均距离的多少倍？
\exsolution
若数密度为 $n$，
典型分子间距定义为
$a=n^{-1/3}$。
硬球分子自由程
$\lambda=1/(\sqrt2\pi d^2n)$，
因此比值
\[
\boxed{\frac{\lambda}{a}
=\frac{n^{-2/3}}{\sqrt2\pi d^2}}.
\]
标准状况
$n=p_0/(k_BT_0)
\simeq2.69\times10^{25}\,\mathrm{m^{-3}}$。
取有效氮分子直径
$d\simeq0.37\,\mathrm{nm}$，
有 $a\simeq3.34\,\mathrm{nm}$、
$\lambda\simeq60\,\mathrm{nm}$，
于是
\[
\boxed{\lambda/a\simeq18}.
\]
\excheck 平均自由程与分子间距不是同一个物理量；
气体虽有相邻分子，
碰撞还取决于分子的硬球碰撞截面。

\thermotransportproblem{225}{88}
一气体在标准状况下，
范德瓦耳斯排除体积参数为
$b=40\,\mathrm{mL/mol}$。
求气体分子的平均自由程。
\exsolution
硬球分子的范德瓦耳斯参数与有效直径 $d$
满足
\[
b=4N_A\left(\frac{\pi d^3}{6}\right)
=\frac{2\pi}3N_A d^3.
\]
因此
\[
d=\left(\frac{3b}{2\pi N_A}\right)^{1/3}.
\]
由硬球平均自由程
$\lambda=k_BT_0/(\sqrt2\pi d^2p_0)$，
代入 $d$：
\[
\boxed{\lambda=
\frac{k_BT_0}{\sqrt2\pi p_0}
\left(\frac{2\pi N_A}{3b}\right)^{2/3}}.
\]
用
$b=40\,\mathrm{mL/mol}=4.0\times10^{-5}\,\mathrm{m^3/mol}$，
$T_0\simeq273.15\,\mathrm K$，
\[
\boxed{\lambda\simeq8.4\times10^{-8}\,\mathrm m
=84\,\mathrm{nm}}.
\]
\excheck 需注意 $b$ 是\emph{每摩尔}的排除体积，
不是单个分子的几何体积。

\thermotransportproblem{226}{88}
声波在标准状况的氮气中传播。
求其波长与氮分子平均自由程相等时的频率。
\exsolution
声速在理想气体中为
\[
c_s=\sqrt{\frac{\gamma RT}{M}},
\quad \gamma\simeq\frac75.
\]
波长等于自由程 $\lambda$ 时
$f=c_s/\lambda$。
由 $\lambda=k_BT/(\sqrt2\pi d^2p)$
和 $R=N_Ak_B$：
\[
\boxed{f=
\frac{\sqrt2\pi d^2p_0}{k_BT_0}
\sqrt{\frac{\gamma RT_0}{M}}
=\frac{\pi d^2p_0N_A}{RT_0}
\sqrt{\frac{2\gamma RT_0}{M}}}.
\]
取 $d\simeq0.37\,\mathrm{nm}$、
$T_0\simeq273\,\mathrm K$：
\[
\boxed{f\simeq5.5\,\mathrm{GHz}}.
\]
\excheck 当波长与分子自由程同量级时，
连续介质的通常声学近似接近失效；
数值是把宏观声速外推到该尺度得到的形式估计。

\thermotransportproblem{227}{88}
温度 $0^\circ\mathrm C$ 的氧气
密封在特征尺度
$\ell=10\,\mathrm{mm}$ 的容器中。
求：(a) 分子平均自由程
$\lambda>\ell$ 所需压强范围；
(b) 临界压强处氧气分子数密度与平均间距。
\exsolution
硬球平均自由程
$\lambda=k_BT/(\sqrt2\pi d^2p)$。
令 $\lambda=\ell$，
得到临界压强
\[
\boxed{p_c=\frac{k_BT}{\sqrt2\pi d^2\ell}},
\quad
\lambda>\ell\iff p<p_c.
\]
取氧分子直径
$d\simeq0.35\,\mathrm{nm}$、
$T=273\,\mathrm K$，
\[
\boxed{p_c\simeq0.7\,\mathrm{Pa}}.
\]
临界状态数密度
$n_c=p_c/(k_BT)$，
平均间距
$a=n_c^{-1/3}$，
\[
\boxed{n_c\simeq2\times10^{14}\,
\mathrm{cm^{-3}}},\qquad
\boxed{a\simeq0.2\,\mathrm{\mu m}}.
\]
\excheck 条件为自由程超过\emph{特征几何尺度}，
不是要让分子间距也超过10 mm；
两者数量级不同。

\thermotransportproblem{228}{88}
标准状况下氮气，求：
(a) 每个分子每秒平均碰撞次数；
(b) $1\,\mathrm{cm^3}$ 氮气中每秒
所有分子间的碰撞总次数。
\exsolution
分子数密度
$n=p_0/(k_BT_0)$，
平均速率
$\bar v=\sqrt{8RT_0/(\pi M)}$。
每分子平均碰撞率
\[
\boxed{z_1=\frac{\bar v}{\lambda}
=\sqrt2\pi d^2 n\bar v}.
\]
采用氮气有效硬球直径约
$d\simeq0.34$--$0.37\,\mathrm{nm}$，
得到
$z_1$ 为
$10^{10}\,\mathrm{s^{-1}}$ 量级；
原书采用的碰撞截面取值给
$\boxed{z_1\simeq0.74\times10^{10}\,\mathrm{s^{-1}}}$。
(b) 每次碰撞涉及两个分子，
计算全部碰撞对时须除以2。
体积 $V=1\,\mathrm{cm^3}$ 中分子数为 $nV$：
\[
\boxed{Z=\frac12nVz_1}.
\]
取 $n\simeq2.7\times10^{19}\,\mathrm{cm^{-3}}$：
\[
\boxed{Z\simeq1.0\times10^{29}\,
\mathrm{s^{-1}}}
\quad\text{（每立方厘米）}.
\]
\excheck 碰撞总次数不能简单写成 $nVz_1$，
否则每个碰撞对会被重复计数两次。

\thermotransportproblem{229}{89}
在理想气体的
(a) 等容过程；(b) 等压过程中，
分子平均自由程 $\lambda$
与每分子单位时间碰撞次数 $z$
如何依赖绝对温度 $T$？
\exsolution
硬球模型中
$\lambda=(\sqrt2\pi d^2n)^{-1}$，
平均分子速率 $\bar v\propto\sqrt T$，
碰撞次数
$z=\bar v/\lambda$。
(a) 等容且物质的量不变，
数密度 $n=N/V$ 不变，
所以
\[
\boxed{\lambda=\text{常数},\qquad z\propto T^{1/2}}.
\]
(b) 等压过程中理想气体
$n=p/(k_BT)\propto T^{-1}$，
故
\[
\boxed{\lambda\propto T,\qquad z\propto T^{-1/2}}.
\]
\excheck 不可把分子平均速率误当成 $\propto T$；
麦克斯韦平均速率正比于 $\sqrt T$，
这决定碰撞频率的半整数幂。
% === END 15-transport-220-229.tex ===

% === BEGIN 15-transport-230-239.tex blob 3f4c4402bfd3950f905c595ea45d13644cd30069 ===
% 原书第二篇第2.7节230--239；PDF页98--99，印刷页89--90。
% 本节凡碰撞频率z均指单分子每秒碰撞次数；
% 全系统碰撞事件数率记Z=(1/2)Nz。
\thermotransportproblem{230}{89}
理想气体压强在某过程中增大 $n$ 倍。
若过程分别为
(a) 等容；(b) 等温，
平均自由程 $\lambda$ 与每分子碰撞频率 $z$
各如何变化？
\exsolution
由硬球运动论
$\lambda\propto1/n_{\rm dens}$、
$z\propto n_{\rm dens}\sqrt T$。
(a) 等容时密度不变，
由 $p\propto T$ 得
$T_f/T_i=n$，
故
\[
\boxed{\lambda_f/\lambda_i=1,\qquad
z_f/z_i=\sqrt n}.
\]
(b) 等温时 $T$ 不变，
由 $p\propto n_{\rm dens}$ 得
$n_{{\rm dens},f}/n_{{\rm dens},i}=n$，
因此
\[
\boxed{\lambda_f/\lambda_i=1/n,\qquad
z_f/z_i=n}.
\]
\excheck 本题的 $n$ 是\emph{压强倍率}，
而 $n_{\rm dens}$ 表示分子数密度，
两者不能混为同一物理量。

\thermotransportproblem{231}{89}
由刚性双原子分子组成的理想气体
经历可逆绝热过程。
求平均自由程 $\lambda$
和每分子碰撞频率 $z$ 分别依赖于：
(a) 容积 $V$；(b) 压强 $p$；
(c) 绝对温度 $T$ 的幂律。
\exsolution
刚性双原子气体
$\gamma=7/5$。
分子数密度 $n_{\rm dens}\propto1/V$，
故 $\lambda\propto V$；
而
$z\propto n_{\rm dens}\sqrt T
\propto V^{-1}T^{1/2}$。
绝热关系
$TV^{\gamma-1}=\text{常数}$，
故 $T\propto V^{-2/5}$。
所以
\[
\boxed{\lambda\propto V,\qquad
z\propto V^{-6/5}}.
\]
又
$pV^\gamma=\text{常数}$，
$V\propto p^{-5/7}$，
得到
\[
\boxed{\lambda\propto p^{-5/7},\qquad
z\propto p^{6/7}}.
\]
最后
$V\propto T^{-5/2}$，
因此
\[
\boxed{\lambda\propto T^{-5/2},\qquad
z\propto T^{3}}.
\]
\excheck 最后一项须由
$V^{-6/5}\propto(T^{-5/2})^{-6/5}=T^3$
计算；若简答文本抽取显示 $T^2$，
与同题前两种表示不能同时相容，应依据连续性核验。

\thermotransportproblem{232}{89}
理想气体经历多方过程
$pV^n=\text{常数}$。
求平均自由程 $\lambda$
和每分子碰撞频率 $z$
分别随
(a) $V$；(b) $p$；(c) $T$
的幂律。
\exsolution
仍有
$\lambda\propto V$、
$z\propto V^{-1}\sqrt T$。
多方状态关系给
$T=pV/(\nu R)\propto V^{1-n}$，
所以
\[
\boxed{\lambda\propto V,\qquad
z\propto V^{-(n+1)/2}}.
\]
若 $n\ne0$，
$p\propto V^{-n}$，
$V\propto p^{-1/n}$，
故
\[
\boxed{\lambda\propto p^{-1/n},\qquad
z\propto p^{(n+1)/(2n)}}.
\]
若 $n\ne1$，
$V\propto T^{1/(1-n)}$，
得到
\[
\boxed{\lambda\propto T^{1/(1-n)},\qquad
z\propto T^{(n+1)/[2(n-1)]}}.
\]
\excheck 特例 $n=0$ 为等压过程，
不能以 $p$ 为随过程变化的自变量；
$n=1$ 是等温过程，
不能用 $T$ 的变化幂律描述。

\thermotransportproblem{233}{89}
刚性双原子理想气体经历多方过程。
若分子间碰撞事件的次数率保持恒定，
求过程摩尔热容：
(a) 每单位体积的碰撞事件率不变；
(b) 整团气体的碰撞事件总率不变。
\exsolution
记分子数密度为 $n_d=N/V$，
单分子碰撞频率
$z\propto n_d\sqrt T$。
\emph{每单位体积的碰撞事件率}
$Z_v=(1/2)n_dz
\propto n_d^2\sqrt T
\propto V^{-2}T^{1/2}$。
(a) 若 $Z_v$ 不变，
则 $T\propto V^4$。
由 $T\propto V^{1-n}$，
多方指数为 $n=-3$。
双原子气体 $C_V=5R/2$，
\[
\boxed{C_a=C_V+\frac R{1-n}
=\frac{5R}{2}+\frac R4
=\frac{11R}{4}\simeq23\,\mathrm{J/(mol\,K)}}.
\]
(b) 若总碰撞率
$Z_{\rm tot}=VZ_v\propto V^{-1}\sqrt T$
不变，则
$T\propto V^2$、$n=-1$，
\[
\boxed{C_b=C_V+\frac R{1-(-1)}
=3R\simeq25\,\mathrm{J/(mol\,K)}}.
\]
\excheck \textbf{与书末简答存在定义冲突：}
本版简答(b) 印 $C=7R/2\simeq29\,\mathrm{J/(mol\,K)}$，
但若“碰撞次数”严格指\emph{单位时间内的碰撞事件率}，
则 $Z_{\rm tot}\propto V^{-1}\sqrt T$ 恒定导出 $3R$。
已明确碰撞计数口径；
若底本采用其他时间尺度上的“碰撞次数”，
需进一步查证，暂不把两答案强行等同。

\thermotransportproblem{234}{89}
摩尔质量为 $M$ 的理想气体装在
体积 $V$ 的薄壁容器中，
容器壁维持温度 $T$。
在 $t=0$ 时打开面积很小的孔口 $S$，
气体开始向真空逸出。
若初始数密度为 $n_0$，
求数密度 $n(t)$。
\exsolution
容器壁保持温度，
可视为逸出气体内部分布始终处于
温度 $T$ 的准平衡麦克斯韦分布。
单位面积每秒外逸粒子通量
$\Phi=n\bar v/4$，
其中
$\bar v=\sqrt{8RT/(\pi M)}$。
容器粒子数 $N=nV$，故
\[
\frac{\mathrm dN}{\mathrm dt}
=-S\frac{n\bar v}{4}
=-S\frac{N\bar v}{4V}.
\]
因此
$\mathrm dN/\mathrm dt=-(1/\tau)N$，
时间常数
$\tau=4V/(S\bar v)$。
解得
\[
\boxed{n(t)=n_0e^{-t/\tau}},\qquad
\boxed{\tau=\frac{4V}{S\bar v}}.
\]
\excheck 需是孔径小于分子平均自由程尺度的
\emph{分子逸出}，不能直接套连续流中的喷嘴方程。

\thermotransportproblem{235}{89}
一容器被隔热薄板分成相等两腔。
隔板有两个孔：一个很小、
一个尺寸远大于分子平均自由程。
第二腔始终维持在第一腔温度的
$\eta$ 倍。
当关闭大孔、只让小孔连通后，
第二腔气体数密度变化几倍？
\exsolution
大孔打开时，
两腔可达常规流体力学的共同压强：
$p_1=p_2$。
由 $p_i=n_i k_BT_i$ 和
$T_2=\eta T_1$，
初态有
$n_{2,0}/n_{1,0}=1/\eta$。
关闭大孔而保留\emph{分子流}小孔时，
分子逸出平衡需要两方向通量相等：
\[
n_{1,f}\bar v_1=n_{2,f}\bar v_2.
\]
因为 $\bar v\propto\sqrt T$，
得到
$n_{2,f}/n_{1,f}=1/\sqrt\eta$。
两腔体积相等且总分子数不变，
$n_1+n_2$ 在前后相等，
于是
\[
n_{2,0}=\frac{n_1+n_2}{1+\eta},
\qquad
n_{2,f}=\frac{n_1+n_2}{1+\sqrt\eta}.
\]
故
\[
\boxed{\frac{n_{2,f}}{n_{2,0}}
=\frac{1+\eta}{1+\sqrt\eta}}.
\]
\excheck $\eta>1$ 时该比值大于1，
即高温第二腔分子浓度增加，
这是一种热渗透/分子逸出平衡效应。

\thermotransportproblem{236}{89}
理想气体经过某过程后，
动力黏度系数增大
$\alpha=2.0$ 倍，
扩散系数增大
$\beta=4.0$ 倍。
求压强变化的倍率。
\exsolution
稀薄气体运动论给
\[
\eta_{\rm visc}\propto\sqrt T,\qquad
D\propto\frac{\sqrt T}{n_d},
\]
其中 $n_d$ 是分子数密度。
黏度增至 $\alpha$ 倍意味着
$T_f/T_i=\alpha^2$。
扩散系数增至 $\beta$ 倍：
\[
\beta=\frac{D_f}{D_i}
=\alpha\frac{n_{d,i}}{n_{d,f}},
\]
故
$n_{d,f}/n_{d,i}=\alpha/\beta$。
由 $p=n_dk_BT$，
\[
\boxed{\frac{p_f}{p_i}
=\frac{\alpha}{\beta}\alpha^2
=\frac{\alpha^3}{\beta}
=\frac{2^3}{4}=2}.
\]
\excheck 原书答案行纯文本看似
“$\alpha^2/\beta$”，但从
$T\sim\eta_{\rm visc}^2$ 与 $D\sim\eta/n_d$
的两个关系检查，压力倍率应为
$\alpha^3/\beta$；给定数值恰为2，
与书末数值答案一致。

\thermotransportproblem{237}{89}
如果理想气体体积增大 $n$ 倍，
求扩散系数 $D$ 与动力黏度 $\eta$
的变化倍数：
(a) 等温膨胀；
(b) 等压膨胀。
\exsolution
动理论关系
$D\sim\bar v\lambda$、
$\eta\sim\rho\bar v\lambda$。
对同一种气体
$\bar v\propto T^{1/2}$，
$\lambda\propto1/n_d\propto V$，
而 $\rho\propto1/V$，因此
$\eta\propto T^{1/2}$、
$D\propto V T^{1/2}$。
(a) $T$ 不变时：
\[
\boxed{\frac{D_f}{D_i}=n,\qquad
\frac{\eta_f}{\eta_i}=1}.
\]
(b) $p$ 不变时理想气体
$T\propto V$，所以：
\[
\boxed{\frac{D_f}{D_i}=n^{3/2},\qquad
\frac{\eta_f}{\eta_i}=n^{1/2}}.
\]
\excheck 黏度不是一般地正比于气体密度；
在稀薄气体动理论中，密度和平均自由程
的变化彼此抵消，留下 $\sqrt T$ 因子。

\thermotransportproblem{238}{89}
刚性双原子理想气体可逆绝热压缩，
体积变成原来的 $1/n$，
其中 $n=10$。
求扩散系数 $D$ 与黏度 $\eta$
各增减多少倍。
\exsolution
刚性双原子气体
$\gamma=7/5$，
绝热压缩有
$T_f/T_i=(V_i/V_f)^{\gamma-1}
=n^{2/5}$。
平均速率变化
$\bar v_f/\bar v_i=n^{1/5}$。
平均自由程与体积成正比，
$\lambda_f/\lambda_i=1/n$，
因此
\[
\frac{D_f}{D_i}
=\frac{\bar v_f}{\bar v_i}
\frac{\lambda_f}{\lambda_i}
=n^{1/5-1}=n^{-4/5},
\]
即
\[
\boxed{D\text{ 减小 }n^{4/5}=10^{0.8}
\simeq6.3\text{ 倍}}.
\]
黏度只随 $\sqrt T$ 变化，
\[
\boxed{\eta\text{ 增大 }n^{1/5}
=10^{0.2}\simeq1.6\text{ 倍}}.
\]
\excheck 不可把扩散系数和黏度
都视为同一密度幂律；二者对体积的依赖不同。

\thermotransportproblem{239}{90}
理想气体沿多方过程
$pV^n=\text{常数}$。
如果在该过程中下列系数保持不变，
分别求多方指数 $n$：
(a) 扩散系数 $D$；
(b) 动力黏度 $\eta$；
(c) 热导率 $\kappa$。
\exsolution
多方关系给
$T\propto V^{1-n}$。
稀薄气体中
$D\propto V\sqrt T
\propto V^{(3-n)/2}$。
(a) 要使指数为零：
\[
\boxed{n_D=3}.
\]
动力黏度
$\eta\propto\sqrt T
\propto V^{(1-n)/2}$，
故
\[
\boxed{n_\eta=1}.
\]
对同一种气体、定容质量比热近似常数，
热导率
$\kappa\propto\rho\bar v\lambda c_V
\propto\sqrt T$，
所以
\[
\boxed{n_\kappa=1}.
\]
\excheck $n=1$ 是等温过程，
而 $n=3$ 是非等温的特殊多方过程；
气体的三种输运系数无需同时不变。
% === END 15-transport-230-239.tex ===

% === BEGIN 15-transport-240-245.tex blob c757ba095700a860660555a2d125caadfc32a71d ===
% 原书第二篇2.7节240--245，用户PDF页99--100，印刷页90。
% 240--242使用经典硬球碰撞近似，超稀薄极限单独列模型。
\thermotransportproblem{240}{90}
已知标准状况下氦气的动力黏度系数，
利用分子运动论计算氦原子的有效直径。
\exsolution
在经典硬球近似下，
$\lambda=(\sqrt2\pi d^2n_d)^{-1}$，
而动理论的黏度公式是
$\eta=\rho\bar v\lambda/3$。
令单原子质量为
$m=M/N_A$，
$\rho=n_dm$，
得到
\[
\eta=\frac{n_dm\bar v}
{3\sqrt2\pi d^2 n_d}
=\frac{m\bar v}{3\sqrt2\pi d^2}.
\]
因此
\[
\boxed{d=\sqrt{
\frac{m\bar v}{3\sqrt2\pi\eta}
}},\qquad
\bar v=\sqrt{\frac{8RT_0}{\pi M}}.
\]
标准状况下取氦气黏度
$\eta\simeq(1.8$--$2.0)\times10^{-5}\,
\mathrm{Pa\,s}$、
$M=0.0040\,\mathrm{kg/mol}$，
得到
\[
\boxed{d\simeq1.8\times10^{-10}\,\mathrm m
=0.18\,\mathrm{nm}}.
\]
\excheck 原题没有明列黏度的数值，
需要使用标准状况氦气物性值；
本结果属于所采用硬球碰撞模型的\emph{有效直径}，
不应与量子化学意义上的原子硬边界完全等同。

\thermotransportproblem{241}{90}
标准状况下，
氦气的热导率是氩气热导率的
$8.7$ 倍。
求氩原子与氦原子有效直径的比
$d_{\rm Ar}/d_{\rm He}$。
\exsolution
气体动理论给
$\kappa=(1/3)\rho c_V\bar v\lambda$。
同为单原子气体，
质量比热
$c_V=(3/2)k_B/m$，
而
$\rho=n_dm$、
$\lambda=(\sqrt2\pi d^2n_d)^{-1}$，
所以
\[
\kappa\propto
(n_dm)\frac1m
\sqrt{\frac{T}{m}}\,
\frac1{d^2n_d}
\propto\frac{\sqrt T}{d^2\sqrt m}.
\]
在相同标准温度下，
\[
\frac{\kappa_{\rm He}}{\kappa_{\rm Ar}}
=\frac{d_{\rm Ar}^2\sqrt{m_{\rm Ar}}}
{d_{\rm He}^2\sqrt{m_{\rm He}}}.
\]
解出
\[
\boxed{\frac{d_{\rm Ar}}{d_{\rm He}}
=\sqrt{
\frac{\kappa_{\rm He}}{\kappa_{\rm Ar}}
\sqrt{\frac{m_{\rm He}}{m_{\rm Ar}}}}}.
\]
用 $M_{\rm He}\simeq4\,\mathrm{g/mol}$、
$M_{\rm Ar}\simeq40\,\mathrm{g/mol}$：
\[
\boxed{d_{\rm Ar}/d_{\rm He}\simeq1.66
\approx1.7}.
\]
\excheck 两气体都是单原子气体，
故其热容量在比例中相消；
不可漏掉原子质量的平方根因子。

\thermotransportproblem{242}{90}
标准状况下，
两根同轴长圆筒之间充满氦气。
圆筒平均半径为 $R$，
间隙 $\Delta R\ll R$，
外圆筒以小角速度 $\omega$ 转动，
内筒静止。
求内筒每单位长度所受摩擦力矩。
当 $\Delta R=6.0\,\mathrm{mm}$ 时，
为使力矩减小 $n=10$ 倍，
应该将氦气压强降到什么量级？
\exsolution
在连续层流黏滞近似下，
两壁相对切向速度为
$\Delta u\simeq\omega R$，
薄隙剪切应力
$\tau_w\simeq\eta\Delta u/\Delta R
=\eta\omega R/\Delta R$。
内筒每单位长度侧面积为 $2\pi R$，
故摩擦力矩每单位长度
\[
\boxed{\mathcal N_0\simeq
\tau_w(2\pi R)R
=\frac{2\pi\eta\omega R^3}{\Delta R}}.
\]
大幅降低气压至自由程超过板间隙，
进入超稀薄气体模型，题给
剪切应力
$\tau_{\rm rare}=\rho\bar v\,\Delta u/6$，
因此
\[
\mathcal N_{\rm rare}
=\frac{\pi}{3}\rho\bar v\,\omega R^3.
\]
要求
$\mathcal N_{\rm rare}=\mathcal N_0/n$：
\[
\rho\bar v=\frac{6\eta}{n\Delta R}.
\]
用
$\eta\simeq\rho\bar v\lambda/3$、
$\lambda=k_BT/(\sqrt2\pi d^2p)$
及 $\rho\bar v\propto p$，
消去密度得到压强的量级：
\[
\boxed{p\simeq\frac{\sqrt2\,k_BT}
{\pi d^2n\Delta R}}.
\]
若氦原子有效直径取
$d\simeq0.18\,\mathrm{nm}$、
$T\simeq273\,\mathrm K$、$n=10$，
\[
\boxed{p\ \text{为 }1\,\mathrm{Pa}\text{ 左右}
\ (\text{原书近似值 }0.7\,\mathrm{Pa})}.
\]
\excheck 连续黏性与自由分子流分别适用于
$\lambda\ll\Delta R$、
$\lambda\gg\Delta R$；
在过渡 Knudsen 数区间需要更精细的动力学模型，
这里的压强只作为量级估计。

\thermotransportproblem{243}{90}
一个气体充满同轴长圆筒之间的空间，
内外半径为 $R_1<R_2$。
外筒以角速度 $\omega$ 缓慢转动，内筒静止。
已知内筒每单位轴向长度上的
黏滞摩擦力矩大小为 $\mathcal N_1$。
设半径 $r$ 处剪切应力
$\tau(r)=\eta r\,\mathrm d\Omega(r)/\mathrm dr$，
求气体动力黏度 $\eta$。
\exsolution
稳定轴对称旋转层流中，
通过每个同轴圆柱面的力矩每单位长度
必须相同：
\[
\mathcal N_1=2\pi r^2\tau(r)
=2\pi\eta r^3
\frac{\mathrm d\Omega}{\mathrm dr}.
\]
因此
$\mathrm d\Omega/\mathrm dr=
\mathcal N_1/(2\pi\eta r^3)$。
由边界
$\Omega(R_1)=0$、
$\Omega(R_2)=\omega$：
\[
\omega=\int_{R_1}^{R_2}
\frac{\mathcal N_1}{2\pi\eta r^3}\,\mathrm dr
=\frac{\mathcal N_1}{4\pi\eta}
\left(\frac1{R_1^2}-\frac1{R_2^2}\right).
\]
解得
\[
\boxed{\eta=\frac{\mathcal N_1}{4\pi\omega}
\left(\frac1{R_1^2}-\frac1{R_2^2}\right)}.
\]
\excheck 必须保留半径平方的倒数差，
且 $\mathcal N_1$ 是\emph{单位长度力矩}，
对应量纲为 $\mathrm{N}$。

\thermotransportproblem{244}{90}
两块半径均为 $a$ 的圆盘具有公共轴线，
间距 $h\ll a$，
其中一块静止，
另一块以小角速度 $\omega$ 旋转。
盘间气体黏度为 $\eta$。
求静止圆盘所受黏滞摩擦力矩。
\exsolution
在距中心半径 $r$ 处，
旋转盘的切向速率为 $\omega r$。
窄间隙中的线性速度梯度近似为
$\partial u_\theta/\partial z=\omega r/h$，
剪切应力
\[
\tau(r)=\eta\frac{\omega r}{h}.
\]
静止圆盘上的环带面积
$\mathrm dA=2\pi r\,\mathrm dr$，
相应力矩微元
\[
\mathrm d\mathcal N=
r\tau(r)\,\mathrm dA
=\frac{2\pi\eta\omega}{h}r^3\,\mathrm dr.
\]
积分 $0\le r\le a$：
\[
\boxed{\mathcal N=
\frac{2\pi\eta\omega}{h}
\int_0^a r^3\,\mathrm dr
=\frac{\pi\eta\omega a^4}{2h}}.
\]
\excheck 盘面只有\emph{一个面对间隙中的气体}
承受该剪切力矩，
不能像整个转盘双侧浸液那样再乘2。

\thermotransportproblem{245}{90}
仍是上一题的两块平行圆盘，
但盘间改为超稀薄气体，
摩尔质量为 $M$、
温度 $T$、压强 $p$。
求静止圆盘所受摩擦力矩。
\exsolution
超稀薄气体中自由程远大于板间距，
不再适用连续剪切
$\eta\omega r/h$。
按本节给出的自由分子流近似，
两板相对切向速率为
$\Delta u=\omega r$，
气体对单位面积的剪切应力为
\[
\tau(r)=\frac16\rho\bar v\,\omega r.
\]
半径为 $r$ 的环带面积
$2\pi r\,\mathrm dr$，故
\[
\mathrm d\mathcal N
=r\tau(r)(2\pi r\,\mathrm dr)
=\frac{\pi}3\rho\bar v\omega r^3\,\mathrm dr.
\]
积分：
\[
\mathcal N=\frac{\pi}{12}\rho\bar v\omega a^4.
\]
由理想气体关系
$\rho=pM/(RT)$，
以及
$\bar v=\sqrt{8RT/(\pi M)}$，
得到
\[
\boxed{\mathcal N=
\frac{\omega a^4p}{3}
\sqrt{\frac{\pi M}{2RT}}}.
\]
\excheck 与连续层流结果不同，
本题的力矩不含板间距 $h$，
而随稀薄气体压强 $p$ 线性变化，
符合自由分子输运极限。
% === END 15-transport-240-245.tex ===

% === BEGIN 15-transport-246-257.tex blob f547039e4cbff54a49aa2566b222ca0e9c0e32ef ===
% I. E. Irodov, Problems in General Physics, Part II §2.7, 246--257.
% 对照用户上传PDF页99--100，印刷页90--91；简答PDF页304 (printed p295).
% \kappa 为热导率，\eta 为动力黏度，热流量密度 j_Q 取向外正。
\thermotransportproblem{246}{90}
利用泊肃叶方程，求在一根长为 $\ell$、内半径为 $a$
的细管中，管两端保持气体压强 $p_1,p_2$ 时，
每单位时间通过管道的气体质量 $\dot m$。
气体为摩尔质量 $M$ 的理想气体，温度恒为 $T$、
动力黏度为 $\eta$。
\exsolution
气体可压缩，沿管道压强变化，因此直接把
不可压缩泊肃叶体积流量 $Q$ 乘入口密度并不正确。
在截面局部采用泊肃叶定律
$Q(x)=-(\pi a^4/8\eta)\mathrm dp/\mathrm dx$，
质量流量为
\[
\dot m=\rho(x)Q(x)
=-\frac{\pi a^4M}{8\eta RT}\,
p(x)\frac{\mathrm dp}{\mathrm dx},
\]
其中用到了等温理想气体关系
$\rho=pM/(RT)$。稳态时 $\dot m$ 与位置无关，
积分从 $x=0$ 到 $\ell$ 得
\[
\dot m\ell
=-\frac{\pi a^4M}{8\eta RT}
\int_{p_1}^{p_2}p\,\mathrm dp
=\frac{\pi a^4M}{16\eta RT}(p_1^2-p_2^2).
\]
因此质量流率的大小为
\[
\boxed{\dot m=\frac{\pi a^4 M}{16\eta RT\ell}
|p_1^2-p_2^2|}.
\]
\excheck 极限 $p_1\to p_2$ 时流率趋零；
当压强差远小于平均压强时，
才可将上式近似为平均密度乘普通泊肃叶流量。

\thermotransportproblem{247}{90}
一根侧面完全绝热的复合导热棒，
由长度 $\ell_1,\ell_2$、热导率
$\kappa_1,\kappa_2$ 的两段串接。
两端分别恒温 $T_1,T_2$。
求两段交界处的稳态温度。
\exsolution
两段截面积相同，稳态下热流密度相等：
\[
\kappa_1\frac{T_1-T_*}{\ell_1}
=\kappa_2\frac{T_*-T_2}{\ell_2}.
\]
两边移项，得到
$(\kappa_1/\ell_1+\kappa_2/\ell_2)T_*
=\kappa_1T_1/\ell_1+\kappa_2T_2/\ell_2$，
故
\[
\boxed{T_*=
\frac{\kappa_1T_1/\ell_1+\kappa_2T_2/\ell_2}
{\kappa_1/\ell_1+\kappa_2/\ell_2}}.
\]
\excheck 当 $\kappa_1/\ell_1\gg\kappa_2/\ell_2$ 时
$T_*\to T_1$，符合第一段热阻很小的极限。

\thermotransportproblem{248}{90}
两根截面积相同、长度为 $\ell_1,\ell_2$、
热导率为 $\kappa_1,\kappa_2$ 的棒首尾相接，
侧面绝热。求与之具有相同稳态导热能力、
总长为 $\ell_1+\ell_2$ 的等效均匀棒的热导率。
\exsolution
等截面串联的热阻相加：
\[
\mathcal R_{\rm th}=
\frac{\ell_1}{\kappa_1 S}+\frac{\ell_2}{\kappa_2 S}
=\frac{\ell_1+\ell_2}{\kappa_{\rm eq}S}.
\]
解得
\[
\boxed{\kappa_{\rm eq}=
\frac{\ell_1+\ell_2}
{\ell_1/\kappa_1+\ell_2/\kappa_2}}.
\]
\excheck 是按长度加权的\emph{调和平均}，
而非算术平均；原版书后公式的 PDF 文本抽取
可能遗漏分式横线，应以热阻相加为准。

\thermotransportproblem{249}{91}
长为 $\ell$ 的均匀导热棒侧面绝热，
材料热导率与温度的关系为
$\kappa(T)=\alpha/T$，
其中 $\alpha>0$。
坐标 $x=0,\ell$ 两端分别维持
$T_1,T_2>0$。求 $T(x)$ 与稳态热流密度。
\exsolution
一维稳态傅里叶定律
$j_Q=-\kappa(T)\mathrm dT/\mathrm dx$
给出
\[
j_Q=-\frac{\alpha}{T}\frac{\mathrm dT}{\mathrm dx}
=-\alpha\frac{\mathrm d(\ln T)}{\mathrm dx}.
\]
$j_Q$ 恒定，故 $\ln T$ 随 $x$ 线性变化：
\[
\ln T(x)=\ln T_1+
\frac{x}{\ell}\ln\frac{T_2}{T_1}.
\]
因此
\[
\boxed{T(x)=T_1\left(\frac{T_2}{T_1}\right)^{x/\ell}},
\qquad
\boxed{j_Q=\frac{\alpha}{\ell}\ln\frac{T_1}{T_2}}.
\]
\excheck 若 $T_1>T_2$，热量沿 $+x$ 方向输运，
$j_Q>0$；原书简答以相反方向定义正热流时
会写成 $\ln(T_2/T_1)$，物理结果相同。

\thermotransportproblem{250}{91}
两块金属的总热容分别为 $C_1,C_2$，
以一根长 $\ell$、截面积 $S$、
热导率 $\kappa$ 很小的棒连接；
整个系统与外界绝热，忽略导热棒自身热容。
初始温差为 $\Delta T_0$。
求金属块温差随时间的变化。
\exsolution
记两块瞬时温度 $T_1(t),T_2(t)$，
热功率 $P=(\kappa S/\ell)(T_1-T_2)$。
能量守恒与各块热平衡方程为
\[
C_1\dot T_1=-P,\qquad C_2\dot T_2=P.
\]
两式相减：
\[
\frac{\mathrm d}{\mathrm dt}(T_1-T_2)
=-\frac{\kappa S}{\ell}
\left(\frac1{C_1}+\frac1{C_2}\right)(T_1-T_2).
\]
积分得
\[
\boxed{\Delta T(t)=\Delta T_0e^{-t/\tau}},\quad
\boxed{\tau^{-1}=\frac{\kappa S}{\ell}
\left(\frac1{C_1}+\frac1{C_2}\right)}.
\]
\excheck 两块总体能量恒定，最终温度相同；
$\tau$ 的量纲为秒。

\thermotransportproblem{251}{91}
两块大平行板间距 $\ell$，
分别恒温 $T_1,T_2$；
板间均匀材料的热导率满足
$\kappa(T)\propto\sqrt{T}$。
求稳态温度分布 $T(x)$，
以 $T_1$ 侧为 $x=0$。
\exsolution
\textbf{按原版图像恢复根号：}
题面是 $\kappa\propto\sqrt{T}$，
纯文本抽取会丢失 $\sqrt{\ }$。
写成 $\kappa(T)=A\sqrt T$，
傅里叶定律给
$j_Q=-A\sqrt T\,T'(x)$，
故 $T^{3/2}$ 随 $x$ 线性变化：
\[
\frac{\mathrm d(T^{3/2})}{\mathrm dx}
=-\frac{3j_Q}{2A}=\text{常数}.
\]
按端点条件，
\[
T(x)^{3/2}
=T_1^{3/2}+\frac{x}{\ell}
(T_2^{3/2}-T_1^{3/2}).
\]
故
\[
\boxed{T(x)=T_1
\left[1+\frac{x}{\ell}
\left(\left(\frac{T_2}{T_1}\right)^{3/2}-1\right)
\right]^{2/3}}.
\]
\excheck 端点 $x=0,\ell$ 分别得到 $T_1,T_2$；
若漏掉根号将错误得到平方根型温度解。

\thermotransportproblem{252}{91}
两块大水平板之间为氦气，间距
$\ell=50\,\mathrm{mm}$，
下板温度 $T_1=290\,\mathrm K$，
上板温度 $T_2=330\,\mathrm K$，
气压接近标准大气压。
求稳态热流密度。
\exsolution
在连续气体传热模型中，
单原子气体的热导率
$\kappa\simeq\frac13\rho c_V\bar v\lambda$，
其中 $\bar v=\sqrt{8RT/(\pi M)}$，
$\lambda=(\sqrt2\pi d^2n)^{-1}$、
$\rho c_V=\frac32nk_B$。
代入整理：
\[
\kappa(T)\simeq
\frac{k_B}{\pi d^2}\sqrt{\frac{RT}{\pi M}}
=A\sqrt T.
\]
由傅里叶定律积分，
朝向低温下板的热流密度大小为
\[
|j_Q|=\frac1\ell\int_{T_1}^{T_2}\kappa(T)\,\mathrm dT
=\boxed{\frac{2k_B}{3\pi^{3/2}d^2\ell}
\sqrt{\frac RM}(T_2^{3/2}-T_1^{3/2})}.
\]
按氦气有效碰撞直径
$d\simeq0.20\,\mathrm{nm}$、
$M\simeq0.004\,\mathrm{kg/mol}$，
\[
\boxed{|j_Q|\simeq40\,\mathrm{W/m^2}}.
\]
\excheck 原题未给 $d$，数值依赖采用的
氦气有效硬球碰撞直径；
热流从330 K上板流向290 K下板。

\thermotransportproblem{253}{91}
两大平板间距 $\ell=5.0\,\mathrm{mm}$，
板间为压强 $p=1.0\,\mathrm{Pa}$ 的氦气。
两板温度分别为 $17^\circ\mathrm C$、
$37^\circ\mathrm C$。
求氦原子平均自由程与热流密度。
\exsolution
平均温度约
$T=(290+310)/2\simeq300\,\mathrm K$。
硬球自由程
\[
\lambda=\frac{k_BT}{\sqrt2\pi d^2p}.
\]
取氦有效直径 $d\simeq0.20\,\mathrm{nm}$，
得
\[
\boxed{\lambda\simeq23\,\mathrm{mm}>\ell}.
\]
故气体处于自由分子输运而非连续傅里叶传热区。
采用原书的自由分子通量近似
\[
|j_Q|=\frac{p\bar v(T_2-T_1)}
{6T(\gamma-1)},\qquad
\bar v=\sqrt{\frac{8RT}{\pi M}},
\]
单原子氦气 $\gamma=5/3$，
$M=0.004\,\mathrm{kg/mol}$，得
\[
\boxed{|j_Q|\simeq22\,\mathrm{W/m^2}}.
\]
\excheck 比较 $\lambda$ 与板间距
决定使用哪一种模型；
若将连续介质傅里叶公式外推到
$\lambda>\ell$，不符合本题设定。

\thermotransportproblem{254}{91}
内外半径为 $R_1<R_2$ 的两根同轴无限长
圆筒之间充满均匀导热物质，
其热导率为常数。两筒分别恒温
$T_1,T_2$。求稳态温度分布 $T(r)$。
\exsolution
取长度 $L$ 的柱形等温面，
稳态径向热功率不随半径变化：
\[
P=-2\pi rL\kappa\frac{\mathrm dT}{\mathrm dr}.
\]
因此 $rT'(r)=\text{常数}$，
积分有 $T(r)=A\ln r+B$。
代入 $T(R_1)=T_1$、
$T(R_2)=T_2$ 得
\[
\boxed{T(r)=T_1+
\frac{T_2-T_1}{\ln(R_2/R_1)}
\ln\frac r{R_1}},\qquad R_1\le r\le R_2.
\]
\excheck $T(r)$ 对 $\ln r$ 线性，
不是对半径 $r$ 线性。

\thermotransportproblem{255}{91}
把前题的两同轴圆筒改为
内外半径 $R_1<R_2$ 的同心球壳，
球面分别恒温 $T_1,T_2$。
求稳态温度分布。
\exsolution
半径 $r$ 的球面上的传热功率为
\[
P=-4\pi r^2\kappa\,T'(r)=\text{常数}.
\]
故 $r^2T'(r)$ 恒定，
积分得到 $T(r)=A+B/r$。
代入两边界值，
\[
\boxed{T(r)=T_1+
(T_2-T_1)
\frac{1/R_1-1/r}{1/R_1-1/R_2}},
\qquad R_1\le r\le R_2.
\]
\excheck 分子在 $r=R_1$ 为零，
在 $r=R_2$ 等于分母，端点温度无误。

\thermotransportproblem{256}{91}
截面半径为 $R$ 的均匀长导线，
内部有恒定电流，使每单位体积产生
恒定焦耳热功率 $w$。
导线热导率为 $\kappa$，
其外表面维持温度 $T_0$。
求截面内的稳态温度分布。
\exsolution
轴对称径向稳态导热方程为
\[
\frac1r\frac{\mathrm d}{\mathrm dr}
\left(r\frac{\mathrm dT}{\mathrm dr}\right)
=-\frac w\kappa.
\]
积分一次：
$rT'(r)=-wr^2/(2\kappa)+C$。
由于轴线上温度梯度有限，$C=0$。
再积分：
$T(r)=-wr^2/(4\kappa)+C_2$。
由 $T(R)=T_0$ 得
\[
\boxed{T(r)=T_0+\frac w{4\kappa}(R^2-r^2)}.
\]
\excheck 中心最热，$T(0)-T_0=wR^2/(4\kappa)$；
正则性条件排除了轴线附近的 $\ln r$ 项。

\thermotransportproblem{257}{91}
半径为 $R$ 的均匀实心球，
内部每单位体积均匀产生热功率 $w$。
热导率为 $\kappa$、球表面温度为 $T_0$。
求球内稳态温度分布。
\exsolution
球对称稳态传热方程为
\[
\frac1{r^2}\frac{\mathrm d}{\mathrm dr}
\left(r^2\frac{\mathrm dT}{\mathrm dr}\right)
=-\frac w\kappa.
\]
积分一次
$r^2T'(r)=-wr^3/(3\kappa)+C$；
由中心有限温度梯度得 $C=0$。
积分并施加表面边界 $T(R)=T_0$，
得到
\[
\boxed{T(r)=T_0+\frac w{6\kappa}(R^2-r^2)}.
\]
\excheck 长圆柱（第256题）系数为 $1/4$，
球体为 $1/6$，反映不同径向通量面积。
% === END 15-transport-246-257.tex ===
